\section{Security Model and Formal Definitions}
\label{sec:model}

\subsection{System Participants and Threat Model}

We consider three principals:
\begin{itemize}
  \item \textbf{Operator} $\mathcal{O}$: controls the AI pipeline and the
  transparency log server. May issue redaction requests. \emph{We model
  $\mathcal{O}$ as a potentially malicious party with respect to selective
  deletion.}
  \item \textbf{Data Subject} $\mathcal{S}$: the natural person whose
  decision is recorded. Initiates erasure requests under GDPR Art.~17.
  \item \textbf{Auditor} $\mathcal{A}$: an independent party (regulatory
  authority, third-party auditor) that verifies the statistical properties
  of the log without accessing individual decision content.
\end{itemize}

\noindent The network adversary $\mathcal{ADV}$ may observe all public log
entries and all published proofs, but may not break the underlying
cryptographic primitives (chameleon hash, ZK-SNARK soundness, Merkle
collision resistance).

\subsection{The Selective Deletion Attack}

\begin{definition}[Selective Deletion Attack]
\label{def:sda}
Let $L = (\ell_1, \ldots, \ell_n)$ be a transparency log of AI decisions.
For each entry $\ell_i$ let $g(\ell_i) \in \mathcal{G}$ denote the
protected-group attribute of the subject and $r(\ell_i) \in \{0,1\}$
denote the decision outcome (0 = reject, 1 = approve). Define the
conditioned approval rate
\[
  \rho(L, g) \;=\; \frac{|\{i : r(\ell_i) = 1 \land g(\ell_i) = g\}|}
                       {|\{i : g(\ell_i) = g\}|}.
\]
An operator mounts a Selective Deletion Attack by issuing a set of
redaction requests $R \subseteq [n]$ such that, for some protected group
$g^* \in \mathcal{G}$,
\[
  \rho(L_{\setminus R},\, g^*) \;>\; \rho(L,\, g^*) + \delta
\]
for a non-negligible $\delta$, while publishing a privacy justification
claiming each redaction is a legitimate GDPR Art.~17 erasure.
\end{definition}

\noindent Existing redactable blockchain schemes~\cite{ateniese2017,deuber2019}
are trivially vulnerable to Definition~\ref{def:sda}: they authenticate
that a redaction was authorised by the key-holder, but place no constraint
on the statistical effect of the redaction batch.

\subsection{The NDR Validity Predicate}

\begin{definition}[$\varepsilon$-Statistical Consistency]
\label{def:eps-stat}
A redaction batch $R \subseteq [n]$ over log $L$ is
$\varepsilon$-\emph{statistically consistent} with respect to a fairness
policy $\Pi = \{(g, \varepsilon_g)\}_{g \in \mathcal{G}}$ if, for every
protected group $g \in \mathcal{G}$,
\[
  \bigl|\rho(L_{\setminus R},\, g) - \rho(L,\, g)\bigr| \;\leq\; \varepsilon_g.
\]
\end{definition}

\begin{definition}[Valid Redaction]
\label{def:valid-redaction}
A redaction batch $R$ is \emph{NDR-valid} for log $L$ under policy $\Pi$
if and only if
\[
  \text{Valid}(\text{Redact}(L, R)) \;\iff\;
  \underbrace{\mathsf{Verify}(\pi_{\mathit{receipt}})}_{
    \text{redaction is logged}}
  \;\land\;
  \underbrace{\mathsf{Verify}(\pi_{\mathit{stat}},\, \varepsilon,\, \Pi)}_{
    \text{$\varepsilon$-stat.~consistency holds}}
\]
where $\pi_{\mathit{receipt}}$ is an inclusion proof that the redaction
event was appended to the log, and $\pi_{\mathit{stat}}$ is a ZK-SNARK
over the Merkle-committed log entries proving Definition~\ref{def:eps-stat}.
\end{definition}

\noindent Note the deliberate contrast with Papers~1 and~2 of this series.
In BTV and BTV-Transparency, accountability is enforced at compile time via
linear resource types; the invariant is a type law checked by the Rust
compiler. In NDR, the invariant operates at the cryptographic layer and is
checked at runtime by a SNARK verifier. The two layers are
complementary: the type system ensures no decision escapes without a
persistence record; the SNARK verifier ensures no persistence record can
be selectively erased without statistical justification. Neither mechanism
subsumes the other.

\subsection{Security Properties}

We require the NDR protocol to satisfy the following properties.

\begin{property}[Redaction Soundness]
No computationally bounded operator $\mathcal{O}$ can produce a proof
$\pi_{\mathit{stat}}$ that passes
$\mathsf{Verify}(\pi_{\mathit{stat}}, \varepsilon, \Pi)$ for a redaction
batch $R$ that violates $\varepsilon$-statistical consistency
(Definition~\ref{def:eps-stat}), except with negligible probability.
\end{property}

\begin{property}[Content Indistinguishability]
For any redacted index $i \in R$, the commitment $C_i'$ published after
redaction is computationally indistinguishable from a fresh null
commitment $C_{\bot}$. In particular, the content of the original decision
$\ell_i$ is irrecoverable from $C_i'$.
\end{property}

\begin{property}[Merkle Continuity]
Let $\mathit{rt}$ and $\mathit{rt}'$ denote the Merkle roots before and
after a valid redaction batch. Any party holding consistency proofs for
$\mathit{rt}$ can verify membership of unredacted entries with respect to
$\mathit{rt}'$ without re-downloading the log.
\end{property}

\begin{property}[Operator Non-Frameability]
A data subject $\mathcal{S}$ or auditor $\mathcal{A}$ cannot forge a
redaction receipt or a failing $\pi_{\mathit{stat}}$ to falsely accuse
the operator of a Selective Deletion Attack.
\end{property}

\subsection{Relationship to Existing Notions}

Redaction Soundness (Property~1) is stronger than the
``policy compliance'' notion of Derler et al.~\cite{derler2019} because
it constrains the aggregate statistical effect of a redaction batch, not
merely the authorisation of individual redactions.

Content Indistinguishability (Property~2) follows from the hiding
property of the underlying Pedersen commitment and does not require
new proof machinery.

Merkle Continuity (Property~3) is standard in Certificate
Transparency~\cite{rfc6962} via the existing consistency proof mechanism;
our construction inherits it without modification.

Operator Non-Frameability (Property~4) is a liveness condition not
present in prior work, added because our threat model elevates the
operator to a potentially malicious party.
